\documentclass{article}
\usepackage{graphicx} % Required for inserting images
\usepackage{amsmath}
\usepackage{amssymb}
\usepackage{physics2}
\usephysicsmodule{ab}
\usepackage{derivative}
\DeclareMathOperator{\sgn}{sgn}
\usepackage{mathrsfs}

\newcommand{\lablace}[1]{\mathscr{L}\ab\{#1\}}

\title{ENAS 1940}
\author{David Chen}
\date{Spring 2026}

\begin{document}

\maketitle

\section{Introductory Statements}
An \textbf{ordinary differential equation (ODE)} is an equation relating ordinary derivatives, and a \textbf{partial differential equation (PDE)} relates partial derivatives.
The \textbf{order} of a differential equation is the highest numbered derivative in the equation.\\
We mostly cover \textbf{linear} differential equations in this course, which allow for the \textbf{superposition} (summation/scaling) of solutions. 

A \textbf{homogeneous} equation has no terms that are not a function of the dependent variable.

\section{Method of Integrating Factors}
Shortening $p = p(t)$, $g=g(t)$, $\mu = \mu(t)$
\begin{align*}
    \odv{y}{t} + py &= g\\
    \mu\ab(\odv{y}{t} + py) &= \mu g\\
    \mu\odv{y}{t} + \mu py &= \mu g
\end{align*}
We choose $\mu$ such that $\mu p y = \odv{\mu}{t}y$, thus satisfying the form of the Product Rule. From such a condition, we get that
$$\frac{\odv{\mu}{t}}{\mu} = p$$
$$\frac{\odif{\mu}}{u} = p\,\odif{t}$$
So $$\mu=e^{\int p\odif{t}} + C$$

For a first-order ordinary differential equation, given that $p,g$ are continuous on a domain, there will exist a unique solution to the differential equation on that domain.

\section{Method of Undetermined Coefficients}
The solution set of a nonhomogeneous (second order) differential equation
$$y''+py'+qy=g$$
can be written in the form
$y=c_1y_1(t)+c_2y_2(t)+Y(t)$
where the first two terms are the homogeneous equation's solutions (\textbf{complementary solution}), and the $Y(t)$ is to be determined.

If the form of $g$ is something familiar like a sinusoid or polynomial, the particular solution may take the form of such terms up to a set of scalars that one solves for (treating each distinct term as a basis).

\section{Variation of Parameters}
Given a nonhomogeneous second order differential equation with homogeneous solutions $y_1,y_2$, we choose\footnote{Alternatives to this choice exist, but this is the one used in our course.} to find a solution of the form $y=c_1y_1+c_2y_2+Y$ where for $u_1(t)$ and $u_2(t)$, $Y=u_1y_1 + u_2y_2$ and $$u_1'y_1+u_2'y_2=0$$
Evaluation of the differential equation using this solution ($0$ing out the terms with $y_1$,$y_2$ based on the assumption that they are homogeneous solutions) leads to the second condition that
$$u_1'y_1'+u_2'y_2'=g$$

To solve the system, we use the Wronskian determinant, which in the 2x2 case is
$$W=\begin{vmatrix}
    y_1 & y_2\\ y_1' & y_2'
\end{vmatrix}$$
and obtain a result
$$Y=-y_1\int\frac{y_2g}{W}\,dt + y_2\int\frac{y_1g}{W}\,dt$$
Linearly independent solutions will have nonzero Wronskians.


\section{Euler Equations}
For a second-order ODE
$$x^2y''+\alpha xy' + \beta y = 0$$
letting $y=e^{rt}$ and solving the quadratic for $r$ will yield
\begin{itemize}
    \item Real roots $r_1,r_2$: $y=c_1e^{r_1t} + c_2e^{r_2t}$
    \item Redundant roots $r$: $y=(c_1+c_2t)e^{rt}$
    \item Complex conjugates $r=\lambda\pm\mu i$: $y=e^{\lambda t}(c_1\cos{\mu t} + c_2\sin{\mu t})$ 
\end{itemize}

For higher order equations, similar results follow, with $n$ redundant roots corresponding with $t^{n-1}e^{rt}$ and below. This is proven using the ansatz from d'Alembert of $y=v(t)y_1$.

\section{Series Solution about an Ordinary Point}
For a differential equation
$$P(x)y''+Q(x)y'+R(x)y=0$$
choosing $x_0$ in an interval where $P(x_0)\ne0$, we can divide to yield
$$y''+py'+qy=0$$
The solution will take the form $$\sum_{n=0}^\infty a_n(x-x_0)^n$$
$p,q$ being analytic at $x_0$ results in an ordinary point.

The convergence of such a series can be proven using the ratio test, convergence of $p,q$ is usually sufficient to justify this too.

\section{Series Solutions for Regular Singular Points}
For simplicity, we shift the system's independent variable such that our initial point is always $x=0$.
Regular Singular Points do not have analytic $p(0),q(0)$, but instead $\lim_{x\to0}xp(x)=p_0$ and $\lim_{x\to0}x^2q(x)=q_0$. Applying a method of ``Euler coefficients times power series,'' we assume
$y=x^r\sum_{n=0}^\infty a_nx^n= \sum_{n=0}^\infty a_nx^{r+n}$
The possible values for $r$ are the roots of the indicial equation
$$F(r)=r(r-1)+p_0r+q_0=0$$
and so long as $r_1\ne r_2+n$ for $n\in\mathbb{Z}$, one can easily obtain two linearly independent solutions.

\section{Laplace Transform}
The Laplace Transform is defined as
$$\mathscr{L}\ab\{f(t)\}=F(s)=\int_0^\infty f(t)e^{-st}\,dt$$
\subsection{Properties of the Laplace Transform}
The Laplace Transform has many properties that make it useful for manipulation  of differential equations. One of the most notable of which is linearity (inherited from integration/summation). Linearity allows us to derive tables of useful Laplace Transforms and their properties, some of which are given below.
\subsubsection{$n^{th}$ Derivative Property}
$$\lablace{f^{(n)}(t)}=s^nF(s)-s^{n-1}f(0)-s^{n-2}f'(0)-\cdots-f^{(n-1)}(0)$$
For the $n=1$ case, this reduces to
$$\lablace{f'(t)}=sF(s)-f(0)$$
\subsubsection{Time/Frequency Shift Property}
$$\lablace{u(t-c)f(t-c)}=e^{-cs}F(s)$$
$$\lablace{e^{cs}f(t)}=F(s-c)$$
(assuming $s>c$ in the latter equation)
\subsubsection{Convolution Property}
Given two Laplace Transformed functions $F(s),G(s)$:
$$\lablace{h(t)}=F(s)G(s)$$ implies that $$h(t)=f(t)*g(t)=\int_0^tf(t-\tau)g(\tau)\,d\tau=\int_0^tf(\tau)g(t-\tau)\,d\tau$$

\subsection{Application of the Laplace Transform}
For a 2nd order ODE
$$ay''+by'+cy=g(t)=$$ the solution will have
$$Y(s)=\Phi(s)+\Psi(s)$$
where $\phi(t)$ represents the transient response and $\psi(t)$ represents the forced response, and are obtained with \begin{align*} \Phi(s)=H(s)\ab[\ab(as+b)y_0+ay_0'] && \Psi(s)=H(s)G(s)
\end{align*}
$$H(s)=\frac{1}{as^2+bs+c}$$

\section{Separation of Variables}
\subsection{Heat Conduction}
Consider a partial differential equation
$\alpha^2\pdv[order=2,2]{u}{x}=\pdv{u}{t}$ representing the conduction of heat through a one-dimensional rod of length $L$ with insulated sides ($u(0,t)=u(L,t)=0$) and initial temperature distribution $u(x,0)=f(x)$. This is a separable differential equation where $u(x,t)=X(x)T(t)$, and rearrangement of the differential equation yields
$$\frac{X''}{X}=\frac{1}{\alpha^2}\frac{T'}{T}=C$$
where $C$ must be a fixed constant (since the two expressions are independent of one another).\\
Casework for $C > 0$, $C = 0$, and $C<0$ shows that it is only possible to satisfy the initial conditions with $C=-\ab(\frac{n\pi}{L})^2$ (eigenvalues). The resulting fundamental solutions all take the form $$u_n(x,t)=e^{-\ab(\frac{n\pi\alpha}{L})^2t}\cdot\sin\ab(\frac{n\pi x}{L})$$ and the full solution would be $u(x,t)=\sum_{n=1}^\infty c_nu_n$. To fulfill the condition of $u(x,0)=f(x)$, one needs to fulfill $f(x)=\sum_{n=1}^\infty c_n\sin\ab(\frac{n\pi x}{L})$ which would require the Fourier Series 
$$c_n=\frac{2}{L}\int_0^L f(x)\sin\ab(\frac{n\pi x}{L})\,dx$$

\subsubsection{Fourier Series}
In general, any periodic function (we have chosen to construct a periodic extension of the given function in the above example) can be decomposed as
$$f(x)=\frac{a_0}{2}+\sum_{n=1}^\infty\ab(a_n\cos\ab(\frac{n\pi x}{L}) + b_n\sin\ab(\frac{n\pi x}{L}))$$
\begin{align*}
    a_n=\frac{1}{L}\int_{-L}^L f(x)\cos\ab(\frac{n\pi x}{L})\,dx\\
    b_n=\frac{1}{L}\int_{-L}^L f(x)\sin\ab(\frac{n\pi x}{L})\,dx\\
\end{align*}

\subsubsection{Bar with Heated Sides}
This case is identical to the bar with insulated sides, but with an added steady-state of
$$v(x)=(T_2-T_1)\frac{x}{L}+T_1$$ from the sides being kept at $T_1$ and $T_2$.

\subsection{Wave Equation on a String}
$$\alpha^2\pdv[order=2,2]{u}{x}=\pdv[order=2,2]{u}{t}$$ results in natural frequencies of $\frac{n\pi a}{L}$ and wavelength $\frac{2L}{n}$.

\subsection{Laplace's Equation}
Fundamentally, Laplace's Equation is
$\nabla u = 0$
\subsubsection{Dirichlet Problem for a Rectangle}
For a rectangle that is $a$ wide and $b$ tall, where 3 of the sides are fixed to $0$, but the right-hand side is fixed to $u(a,y)=f(y)$ we have
$$u_n(x,y)=\sinh\ab(\frac{n\pi x}{b})\sin\ab(\frac{n\pi y}{b})$$
$$c_n=\frac{2}{b\sinh\ab(\frac{n\pi a}{b})}\int_0^bf(y)\sin\ab(\frac{n \pi y}{b})\,dy$$

\subsubsection{Dirichlet Problem for a Circle}
For a circle of radius $a$ and boundary $u(a,\theta)=f(\theta)$, the relevant form of Laplace's Equation is
$$u_rr+\frac{1}{r}u_r+\frac{1}{r^2}u_{\theta\theta}=0$$
with the values of $f$ being bounded and periodic.
Fundamental solutions (including a constant term) include 
\begin{align*}
    u_n(r,\theta) = r^n\cos(n\theta) && u_n(r,\theta) = r^n\sin(n\theta)
\end{align*}
with the coefficients all corresponding to the relevant Fourier series.

\section{Sturm-Liouville Problems}
Differential Equations of the form
$$(py')' - qy + \lambda ray=0$$
which we assume to have $p,p',q,r$ continuous, $p,r>0$ for $0\le x\le1$ (regular) and boundary conditions
\begin{align*}
    \alpha_1y(0)+\alpha_2y'(0)=0&&\beta_1y(1)+\beta_2y'(1)=0
\end{align*}

The eigenvalues $\lambda$ will always be real-valued with a single eigenfunction each, and the eigenfunctions $\phi_m,\phi_n$ corresponding to distinct eigenvalues will always be orthogonal
$$\int_0^1 r\phi_m\phi_n\,dx=0$$

We can normalize the eigenvalues by dividing their ``magnitude''
$$\int_0^1r\phi_n^2\,dx$$

And from these eigenfunctions, we form a basis from which other functions can be constructed.

\section{Common Trigonometric Identities}
$$a\cos{x}+b\cos{x}=\sgn{(a)}\sqrt{a^2+b^2}\cos{\left(x+\arctan\left(-\frac{b}{a}\right)\right)}$$
$$\cos{\ab(a+b)} = \cos{a}\cos{b} - \sin{a}\sin{b}$$
$$\sin{\ab(a+b)} = \sin{a}\cos{b} + \cos{a}\sin{b}$$
$$\cos{2\theta} = 2\cos^2\theta-1 =  1 - 2\sin^2\theta$$

$$\int x\sin(ax)\,dx = \frac{\sin(ax)}{a^2}-\frac{x\cos(ax)}{a}+C$$
$$\int x\cos(ax)\,dx = \frac{\cos(ax)}{a^2}+\frac{x\sin(ax)}{a}+C$$

\section{Even and Odd Function Identities}
Let $E$ represent some even function, $O$ represent some odd function.
\begin{align*}
    E \cdot E = E && O \cdot E = E &&
    O \cdot O = E
\end{align*}
\begin{align*}
    \int_{-L}^{+L} E(x)\,dx = 2\int_0^{+L}E(x)\,dx&&
    \int_{-L}^{+L} O(x)\,dx = 0
\end{align*}

\end{document}

